El \textit{Shared Package} es un submódulo de Git integrado en ambos proyectos Unity cuyo propósito es
centralizar el código compartido entre el \textit{World Editor} y el proyecto de juego. 
\vspace{0.25cm}
Contiene los DTOs (\textit{Data Transfer Objects}) que definen las estructuras de datos comunes, como los modelos de
\textit{Placeables}, \textit{Creatables} entre otras. Por otro lado, se encapsulan servicios de comunicación
con los endpoints REST del Core Service, permitiendo que tanto el editor como el cliente del juego
consuman la misma capa de acceso a la API sin duplicar su implementación. Al estar versionado como
submódulo independiente, cualquier modificación en las estructuras compartidas se refleja de forma
consistente en ambos proyectos, evitando duplicación de código y reduciendo el riesgo de inconsistencias
entre el editor y el juego.

\subsubsection{Integración como Package de Unity}

Unity cuenta con un sistema de paquetes (\textit{Package Manager}) que permite incorporar librerías y
módulos externos directamente desde un repositorio de Git. El \textit{Shared Package} se registra en
cada proyecto mediante el archivo \texttt{manifest.json} del Package Manager, apuntando a la URL del
repositorio Git correspondiente. Esto permite que Unity resuelva y cargue el paquete automáticamente,
tratándolo como cualquier otra dependencia del proyecto sin necesidad de copiar el código manualmente
entre proyectos.

\subsubsection{Mantenimiento}
Dado que el \textit{Shared Package} es una dependencia compartida entre ambos proyectos, cualquier
modificación en él debe evaluarse con cuidado. Al desarrollar una nueva funcionalidad en el \textit{World Editor}
o en el proyecto de juego, si dicha funcionalidad requiere modificar el \textit{Shared Package}, es
responsabilidad del desarrollador verificar que los cambios introducidos no generen efectos secundarios
ni rompan la compilación del otro proyecto. Esto implica validar ambos proyectos Unity ante cualquier
modificación en los DTOs o servicios compartidos antes de integrar los cambios.